% Drawing 74/04-15 -- folder 74, batch 04, archive page 77 (PDF page 78).
% Transcription: le long du bord droit de la feuille, quatre petits dessins : en haut, un hexagone dont trois côtés sont repassés ; plus bas, une étoile formée de deux triangles et un petit hexagone aux sommets marqués ; au-dessous, au-dessus du petit diagramme des $H$, un hexagone dont deux sommets sont marqués d'un gros point et trois côtés repassés.
% Kind: four small hexagons (heavy sides, star of two triangles, marked vertices). Method 2 (schematic redraw from the facsimile). Status: partial.
% Notes: Along the right edge of the sheet, four small figures drawn as regular figures at the places and sizes of the page: at the top a hexagon with three sides gone over (left, lower left, upper right); lower, a star of two triangles and a small hexagon with its six vertices marked; lower still, on top of the small diagram of the H's, a hexagon with heavy points on its upper left and upper right vertices and three alternate sides gone over. The writing between them (the numbers 3, 2, 2, 6, 6, 12 of the list, struck words) is text and is not drawn; the arrow diagram of the H's under the last hexagon is a relation diagram and is not drawn. Still differs: the figures are regularised; which sides are heavy is read off a 200 dpi render; small arrow-like ticks on the sides of the top hexagon and of the small hexagon (orientations) are left out. Id 04-15: existing ids are not renumbered. Redrawn 2026-11-12 (issue 128): to check again.
% Coordinates: 1 unit = 1 mm = 1/100 of the crop width (bbox in page fractions: [0.85, 0.085, 0.97, 0.44]).
% Redrawn by Claude (claude-fable-5-2), issue #128, 2026-11-12; not checked by a human; not compiled (no TeX in the session).
% Rights: reproduces part of the work -- behind RENDER_TRANSCRIPTS, never published by itself.
\documentclass[tikz,border=3pt]{standalone}
\usepackage{amsmath,amssymb}
\usetikzlibrary{patterns,arrows.meta}
\tikzset{pencil/.style={black!45}, struck/.style={black!35}, mpt/.style={circle,draw,inner sep=1.3pt}, pt/.style={circle,fill,inner sep=0.9pt}, opt/.style={circle,draw,fill=white,inner sep=1.1pt}}
\begin{document}
\begin{tikzpicture}[x=1mm,y=1mm,line cap=round,line join=round,font=\small]
\draw (33.5,-9) -- (19.6,-17);
\draw[very thick] (19.6,-17) -- (19.6,-33);
\draw[very thick] (19.6,-33) -- (33.5,-41);
\draw (33.5,-41) -- (47.4,-33);
\draw (47.4,-33) -- (47.4,-17);
\draw[very thick] (47.4,-17) -- (33.5,-9);
\draw (60,-322) -- (50.5,-338.5) -- (69.5,-338.5) -- cycle;
\draw (50.5,-327.5) -- (60,-344) -- (69.5,-327.5) -- cycle;
\draw (84,-353) -- (75.3,-358) -- (75.3,-368) -- (84,-373) -- (92.7,-368) -- (92.7,-358) -- cycle;
\node[pt,inner sep=0.6pt] at (84,-353) {};
\node[pt,inner sep=0.6pt] at (75.3,-358) {};
\node[pt,inner sep=0.6pt] at (75.3,-368) {};
\node[pt,inner sep=0.6pt] at (84,-373) {};
\node[pt,inner sep=0.6pt] at (92.7,-368) {};
\node[pt,inner sep=0.6pt] at (92.7,-358) {};
\draw[very thick] (33,-379) -- (17.4,-388);
\draw (17.4,-388) -- (17.4,-406);
\draw[very thick] (17.4,-406) -- (33,-415);
\draw (33,-415) -- (48.6,-406);
\draw[very thick] (48.6,-406) -- (48.6,-388);
\draw (48.6,-388) -- (33,-379);
\node[pt,inner sep=1.4pt] at (17.4,-388) {};
\node[pt,inner sep=1.4pt] at (48.6,-388) {};
\end{tikzpicture}
\end{document}
